% =====================================================================
% RASCUNHO -- seção "Filtros de Covariância"
% ---------------------------------------------------------------------
% Onde entra: depois do parágrafo "Covariância é uma medida estatística..."
% (primeiro parágrafo da seção Filtros de Covariância do capítulo).
%
% Chaves de citação NOVAS (ainda não estão no biblio.bib):
%   zhao-lmmse, tse-viswanath, haykin-sistemas, tomasin-interference
% As entradas estão em rascunho_referencias.bib, com os campos a conferir.
%
% Cada afirmação tem um comentário "% FONTE:" com a página ou a linha
% exata, para conferência. O exemplo de duas antenas (tópicos 4 a 7) é
% uma ilustração própria, não vem de nenhuma das fontes.
%
% Cuidado de posição: esta seção vem ANTES da eq. (5) e da construção de
% R. Por isso o texto não usa H barra, beta nem sigma^2.
% =====================================================================


% ---------------------------------------------------------------------
% 1. O equalizador e o critério do erro quadrático médio
% ---------------------------------------------------------------------

No \textit{uplink}, cada antena da estação base recebe uma combinação dos sinais transmitidos pelos UEs, somada ao ruído. Esse sinal recebido pode ser escrito como $\mathbf{y}=\mathbf{H}\mathbf{s}+\mathbf{n}$, em que $\mathbf{H}$ é a matriz de canal entre as $M$ antenas e os $K$ UEs alvo, $\mathbf{s}$ é o vetor de símbolos transmitidos e $\mathbf{n}$ é o ruído \cite{zhao-lmmse}.
% FONTE: artigo base, eq. (1), paper.md:56-60.
O equalizador aplica uma matriz $\mathbf{W}$ ao sinal recebido para obter a estimativa dos símbolos, $\hat{\mathbf{s}}=\mathbf{W}\mathbf{y}$ \cite{zhao-lmmse}.
% FONTE: artigo base, paper.md:79.
Cada linha de $\mathbf{W}$ possui um peso para cada antena. Assim, a estimativa de cada símbolo é uma combinação dos sinais de todas as antenas.
% FONTE: consequência direta das dimensões de W (K x M) na eq. (3), paper.md:75.

No equalizador LMMSE, os pesos são escolhidos para que o erro quadrático médio entre a estimativa e os símbolos transmitidos seja o menor possível \cite{zhao-lmmse}:

$$
\min_{\mathbf{W}}\ \mathbb{E}\left[\|\mathbf{W}\mathbf{y}-\mathbf{s}\|_2^2\right].
$$
% FONTE: artigo base, eq. (2), paper.md:63-67.

A solução desse problema depende da matriz de covariância do ruído, $\mathbf{R}\triangleq\mathbb{E}[\mathbf{n}\mathbf{n}^H]$ \cite{zhao-lmmse}. Para entender por que essa matriz aparece na solução, é preciso ver o que ela diz sobre o ruído de cada antena.
% FONTE: definição de R em paper.md:60; R^{-1} aparece na eq. (3), paper.md:75.


% ---------------------------------------------------------------------
% 2. Assinatura espacial
% ---------------------------------------------------------------------

Um sinal transmitido chega a todas as antenas da estação base. Em cada antena, ele é multiplicado pelo ganho do canal entre o transmissor e aquela antena. Para um único transmissor, o sinal recebido pode ser escrito como $\mathbf{y}=\mathbf{h}s+\mathbf{n}$, em que o vetor $\mathbf{h}$ reúne os ganhos do canal em cada antena \cite[p.~297]{tse-viswanath}. Esse vetor é chamado de assinatura espacial do sinal \cite[p.~297]{tse-viswanath}. Ele descreve o formato com que aquele sinal aparece no conjunto de antenas.
% FONTE: Tse e Viswanath, seção 7.2, eq. (7.18), p. 297 ("spatial signature").


% ---------------------------------------------------------------------
% 3. Ruído independente e ruído correlacionado entre antenas
% ---------------------------------------------------------------------

Quando a covariância entre o ruído de duas antenas é zero, esses ruídos são descorrelacionados \cite[p.~202]{haykin-sistemas}. Se o ruído for gaussiano, eles também são independentes \cite[p.~222]{haykin-sistemas}. Para um vetor gaussiano, isso acontece quando a sua matriz de covariância é diagonal \cite[p.~499]{tse-viswanath}. Nesse caso, conhecer o ruído de uma antena não muda o que se sabe sobre o ruído da outra \cite[p.~196]{haykin-sistemas}. O ruído térmico costuma ser modelado dessa forma, com matriz de covariância diagonal \cite{zhao-lmmse}.
% FONTE: livro de probabilidade, p. 202 (descorrelacionadas se e somente se a covariância é zero);
%        p. 222, Propriedade 4 (gaussianas descorrelacionadas são independentes);
%        p. 196 (independência: conhecer X não afeta a distribuição de Y).
%        Tse e Viswanath, Apêndice A.1.2, item (b), p. 499 (componentes independentes
%        exatamente quando a matriz de covariância é diagonal).
%        Artigo base, paper.md:33 (AWGN com matriz de covariância diagonal).

Já um sinal interferente, transmitido por um UE que não é alvo da estação base, também chega a todas as antenas com a sua própria assinatura espacial. Assim, a mesma interferência aparece ao mesmo tempo em várias antenas, com ganhos diferentes. Essa relação não vem de uma antena interferir na outra, e sim do mesmo sinal chegar a todas elas. Dessa forma, a matriz de covariância do ruído passa a ter elementos fora da diagonal \cite[p.~360]{tse-viswanath}, e o ruído das antenas deixa de ser independente \cite[p.~202]{haykin-sistemas}.
% FONTE: Tse e Viswanath, eq. (8.64)-(8.65), p. 359-360: a covariância do ruído mais
%        interferência é N0 I + soma de P_i h_i h_i^*.
%        NOTA: no livro, a interferência vem dos outros fluxos de um enlace MIMO, não de
%        UEs não-alvo. A álgebra é a mesma; o cenário físico é outro.
%        Livro de probabilidade, p. 202: independentes implica descorrelacionadas; logo,
%        covariância diferente de zero implica que não são independentes.


% ---------------------------------------------------------------------
% 4. Exemplo com duas antenas (ILUSTRAÇÃO PRÓPRIA)
% ---------------------------------------------------------------------

Para ilustrar essa ideia, considere uma estação base com duas antenas, um UE alvo e um UE interferente. O UE alvo transmite o símbolo $s$ com assinatura espacial $\mathbf{h}=[1\;\;1]^T$. O interferente transmite o símbolo $\overline{s}$ com assinatura espacial $\overline{\mathbf{h}}=[1\;\;2]^T$. Cada antena também recebe o seu próprio ruído térmico, $w_1$ e $w_2$. Assim, os sinais recebidos nas duas antenas são

$$
y_1 = s + \overline{s} + w_1, \qquad y_2 = s + 2\overline{s} + w_2.
$$
% NOTA: w_1 e w_2 são o ruído térmico. No resto do capítulo ele é o vetor n barra;
%       decidir se unifica a notação aqui.

Como o ganho do interferente na antena 2 é o dobro do ganho na antena 1, a interferência na antena 2 é sempre o dobro da interferência na antena 1. A Tabela~\ref{tab:interferencia_duas_antenas} mostra essa relação em três instantes.

\begin{table}[htbp]
\caption{Interferência recebida nas duas antenas do exemplo, em três instantes.}
\label{tab:interferencia_duas_antenas}
\centering
\begin{tabular}{cccc}
\hline
Instante & $\overline{s}$ & Interferência na antena 1 & Interferência na antena 2 \\
\hline
1 & $3$ & $3$ & $6$ \\
2 & $-1$ & $-1$ & $-2$ \\
3 & $0{,}5$ & $0{,}5$ & $1$ \\
\hline
\end{tabular}
\legend{Fonte: elaborado pela autora.}
\end{table}

Considerando potência 1 para o interferente e potência $0{,}1$ para o ruído térmico de cada antena, a matriz de covariância do ruído desse exemplo é

$$
\mathbf{R}=
\underbrace{\begin{bmatrix} 1\cdot 1 & 1\cdot 2 \\ 2\cdot 1 & 2\cdot 2 \end{bmatrix}}_{\text{interferente}}
+
\underbrace{\begin{bmatrix} 0{,}1 & 0 \\ 0 & 0{,}1 \end{bmatrix}}_{\text{ruído térmico}}
=
\begin{bmatrix} 1{,}1 & 2 \\ 2 & 4{,}1 \end{bmatrix}.
$$
% FONTE da forma "potência x h h^H + ruído térmico x I": Tse e Viswanath, eq. (8.65), p. 360.
% Números: próprios (conferidos em Python).

Os elementos da diagonal são a potência do ruído em cada antena. Na antena 1, a interferência contribui com 1 e o ruído térmico com $0{,}1$. Na antena 2, a interferência contribui com 4, porque o seu ganho é 2. Já os elementos fora da diagonal valem 2 e vêm apenas do interferente, porque o ruído térmico das duas antenas é independente. Esses elementos mostram o quanto o ruído das duas antenas varia em conjunto.


% ---------------------------------------------------------------------
% 5. Primeiro extremo: anular a interferência
% ---------------------------------------------------------------------

Como a interferência na antena 2 é sempre o dobro da interferência na antena 1, o receptor pode eliminá-la. Para isso, basta combinar as duas antenas com os pesos 2 e $-1$:

$$
2y_1 - y_2 = (2-1)\,s + (2-2)\,\overline{s} + (2w_1 - w_2) = s + (2w_1 - w_2).
$$

A interferência desaparece porque os pesos anulam a assinatura do interferente, já que $2\cdot 1 - 1\cdot 2 = 0$. O sinal do UE alvo permanece porque a sua assinatura é diferente, e $2\cdot 1 - 1\cdot 1 = 1$. Esse receptor projeta o sinal recebido na direção ortogonal à assinatura do interferente, e é chamado de \textit{decorrelator} ou \textit{interference nuller} \cite[p.~81--82]{tse-viswanath}.
% FONTE: Tse e Viswanath, seção 3.3, eq. (3.99) e Fig. 3.13, p. 81 (projeção ortogonal
%        a h_2 elimina a interferência da antena 2); p. 82 ("interference nuller,
%        zero-forcing receiver or the decorrelator"; duas antenas podem anular um
%        usuário interferente).
%        NOTA: o livro diz "projetar"; "combinar com pesos" é a mesma operação escrita
%        antena por antena.


% ---------------------------------------------------------------------
% 6. O custo de cada extremo
% ---------------------------------------------------------------------

Anular a interferência tem um custo. No exemplo, o ruído térmico passa a ser $2w_1 - w_2$, e a sua potência sobe para $4\cdot 0{,}1 + 0{,}1 = 0{,}5$. Além disso, se a assinatura do UE alvo fosse parecida com a do interferente, o sinal desejado também seria quase anulado \cite[p.~81]{tse-viswanath}. Dessa forma, esse receptor elimina a interferência sem considerar quanto do sinal desejado é perdido \cite[p.~358]{tse-viswanath}.
% FONTE: Tse e Viswanath, p. 81 (se a componente de h_1 perpendicular a h_2 é pequena,
%        a recepção piora: "the price paid for nulling out the interference");
%        p. 358 ("without any regard to how much energy of the stream of interest is lost").

No outro extremo, o receptor pode somar as antenas na direção da assinatura do UE alvo. Essa é a melhor escolha quando o ruído é branco \cite[p.~359]{tse-viswanath}, e preserva a energia do sinal desejado \cite[p.~358]{tse-viswanath}. Porém, no exemplo, com os pesos 1 e 1 a interferência não é reduzida:

$$
y_1 + y_2 = 2s + 3\overline{s} + (w_1 + w_2).
$$
% FONTE: Tse e Viswanath, p. 359 ("If the noise is white, we know that it is optimal to
%        project y onto the direction along h"); p. 358 (matched filter preserva a
%        energia do fluxo de interesse).


% ---------------------------------------------------------------------
% 7. O equilíbrio do equalizador LMMSE
% ---------------------------------------------------------------------

O receptor LMMSE busca o equilíbrio entre esses dois extremos, entre combater a interferência e combater o ruído térmico \cite[p.~358]{tse-viswanath}. Para isso, ele calcula os pesos a partir da inversa da matriz de covariância do ruído. Para um único sinal com assinatura $\mathbf{h}$, os pesos são $\mathbf{v}=\mathbf{R}^{-1}\mathbf{h}$ \cite[p.~359]{tse-viswanath}. Quando a interferência é fraca, esse receptor se aproxima da soma na direção do sinal. Quando a interferência é forte, ele se aproxima do receptor que anula a interferência \cite[p.~360]{tse-viswanath}.
% FONTE: Tse e Viswanath, p. 358 (receptor linear que equilibra interferência e ruído);
%        eq. (8.61), p. 359 (v_mmse = K_z^{-1} h, que também minimiza o erro quadrático
%        médio); p. 360 (baixa potência: matched filter; alta: decorrelator).
%        NOTA: no livro a matriz se chama K_z; aqui foi chamada de R para seguir o artigo.

No exemplo, $\mathbf{R}^{-1}\mathbf{h}$ é proporcional aos pesos $2{,}1$ e $-0{,}9$. Esses pesos são próximos de 2 e $-1$, porque a interferência é forte. A Tabela~\ref{tab:comparacao_receptores} compara os três receptores pela razão entre a potência do sinal e a soma das potências de interferência e de ruído, chamada de SINR \cite[p.~358]{tse-viswanath}.
% FONTE dos números: próprios (conferidos em Python). O valor de SINR do LMMSE
%        coincide com a expressão h^* R^{-1} h, eq. (8.62), p. 359.

\begin{table}[htbp]
\caption{Comparação dos três receptores no exemplo de duas antenas.}
\label{tab:comparacao_receptores}
\centering
\begin{tabular}{lccccc}
\hline
Receptor & Pesos & Sinal & Interferência & Ruído térmico & SINR \\
\hline
Soma das antenas & $1$ e $1$ & $4$ & $9$ & $0{,}2$ & $0{,}43$ \\
Anulação da interferência & $2$ e $-1$ & $1$ & $0$ & $0{,}5$ & $2{,}00$ \\
LMMSE & $2{,}1$ e $-0{,}9$ & $1{,}44$ & $0{,}09$ & $0{,}52$ & $2{,}35$ \\
\hline
\end{tabular}
\legend{Fonte: elaborado pela autora.}
\end{table}

O receptor LMMSE deixa passar uma pequena parte da interferência, mas mantém mais sinal que o receptor que anula a interferência. Assim, a sua SINR é maior que a dos dois extremos.


% ---------------------------------------------------------------------
% 8. Onde a matriz de covariância entra
% ---------------------------------------------------------------------

O receptor não conhece a assinatura do interferente nem o símbolo que ele transmite. Mesmo assim, ele consegue calcular os pesos, porque a matriz de covariância do ruído guarda o formato dessa assinatura. No exemplo, na parte da matriz causada pelo interferente, a segunda linha é o dobro da primeira, o mesmo formato de $\overline{\mathbf{h}}=[1\;\;2]^T$.
% FONTE: leitura própria do exemplo; a forma h h^H vem de Tse e Viswanath, eq. (8.65).

Com vários UEs alvo, o equalizador LMMSE é dado por

$$
\mathbf{W}=\left(\mathbf{H}^H\mathbf{R}^{-1}\mathbf{H}+\frac{1}{E_s}\mathbf{I}\right)^{-1}\mathbf{H}^H\mathbf{R}^{-1},
$$

em que $E_s$ é a energia média por símbolo \cite{zhao-lmmse}. Essa expressão não é igual à do exemplo, porque considera ao mesmo tempo os $K$ UEs alvo e a energia dos símbolos. Porém, ela também usa a inversa da matriz de covariância do ruído para ponderar as antenas.
% FONTE: artigo base, eq. (3), paper.md:71-79 (cita Kay [28] como origem).
%        NOTA: v = R^{-1} h (Tse) é um vetor para um único sinal; W da eq. (3) é uma
%        matriz K x M. Não apresentar as duas como a mesma fórmula.

Por isso, a matriz de covariância usada no equalizador precisa representar bem o ruído. A estimativa da correlação da interferência é usada para processar os sinais recebidos e diminuir o efeito da interferência \cite{tomasin-interference}. Além disso, erros nessa estimativa afetam a taxa de dados de forma não linear \cite{tomasin-interference}.
% FONTE: Tomasin et al. [22], pilot-estimation/Estimation*.md:68 ("used to process the
%        forthcoming received signals ... and mitigate the effects of interference");
%        :624 ("the impact of correlation matrix estimation errors is not linear on the
%        data rates").
% LIGAÇÃO: remeter a esta seção em "O que R representa", "O que cada amostra contribui"
%          e "Quantidade de amostras".


% ---------------------------------------------------------------------
% 9. Fora do escopo: métodos baseados no filtro de Kalman
% ---------------------------------------------------------------------

Trabalhos anteriores também trataram receptores MMSE com ruído colorido em sistemas distribuídos, usando um método baseado no filtro de Kalman \cite{zhao-lmmse}. Esses trabalhos consideram que o ruído colorido vem do erro de estimativa do canal. Além disso, assumem que a matriz de covariância é bloco-diagonal, e que cada nó conhece apenas a sua parte local \cite{zhao-lmmse}. Como a covariância considerada neste trabalho possui blocos fora da diagonal, esse método não entra no escopo.
% FONTE: artigo base, Remark 1, paper.md:142 (referências [23] e [24] do artigo).
